\chapter{Linear Programming background}\label{chap:lp-background}
This chapter introduces the linear programming concepts needed to develop a
compositional language for optimization. Readers already familiar with linear
programming may proceed directly to Section~\ref{sec:lp-notions}, which specifies
the definitions and conventions used in this work.

\section{The structure of linear optimization}

Optimization concerns the search for an object that minimizes or maximizes a
specified criterion over a given set. Linear programs form an important class
of optimization problems, with applications in engineering, economics, and
other areas of applied mathematics.

An optimization problem has two components: constraints and an objective.
The constraints define the \emph{feasible region}, the set of admissible choices.
For example, when travelling between cities, admissible routes may be restricted
by a fuel budget, fuel consumption per kilometre, or the money available for
refuelling. A route that violates these restrictions is infeasible.

To obtain a continuous linear program, we model a divisible amount of traffic
rather than the choice of one indivisible car route. Consider three cities
$A,B,C$, with roads $A\to B$, $B\to C$, and $A\to C$. Let
$x_{AB},x_{BC},x_{AC}\ge0$ be the amounts transported along these roads,
let $p\ge0$ be the amount to transport from $A$ to $C$, and let $F\ge0$
be the total fuel budget. Suppose each road consumes one unit of fuel per
unit transported. The constraints are
\[
\begin{aligned}
x_{AB}+x_{AC}&=p,\\
x_{BC}&=x_{AB},\\
x_{AB}+x_{BC}+x_{AC}&\le F,\\
x_{AB},x_{BC},x_{AC}&\ge0.
\end{aligned}
\]
The first two equations express conservation of traffic; the inequality
expresses the fuel restriction. Requiring a single unsplit car route instead
would require a discrete model, which is not the continuous LP considered here.

Several feasible transportation plans may reach the same destination. The
objective provides a criterion for comparing them, such as monetary cost or
travel time. An optimal solution is a feasible choice that achieves the best
objective value, when such a choice exists.

For the continuous transportation model, suppose the monetary costs per unit
on the three roads are $1$, $2$, and $4$, respectively. The objective is
\[
J(x_{AB},x_{BC},x_{AC})=x_{AB}+2x_{BC}+4x_{AC}.
\]
The route through $B$ costs less per unit than the direct road, but consumes
more fuel. Thus the objective and the fuel constraint describe different
aspects of the same choice.

Together, the constraints and objective define the optimization problem:
minimize the transportation cost over the feasible flows.

Combining the objective and constraints gives the value function
\[
V(p,F)=\inf_{x_{AB},x_{BC},x_{AC}\ge0}
\left\{x_{AB}+2x_{BC}+4x_{AC}\;\middle|\;
\begin{aligned}
x_{AB}+x_{AC}&=p,\\
x_{BC}&=x_{AB},\\
x_{AB}+x_{BC}+x_{AC}&\le F
\end{aligned}\right\}.
\]
Here $p$ and $F$ are parameters, while the three road flows are decision
variables. For example, $V(1,1)=4$, since only the direct road is feasible,
whereas $V(1,2)=3$, since the cheaper route through $B$ becomes feasible.
If $p\ge0$ and $F<p$, the problem is infeasible and its value is $+\infty$.

A linear program has linear equality and inequality constraints and a linear
objective, possibly with a constant offset. The feasible region is a polyhedron: an intersection of finitely many closed half-spaces. A nonempty feasible region defined only by linear equalities is an affine subspace, but inequalities generally make it a polyhedron rather than an affine subspace.

Now consider a transportation problem planned in two stages. The first stage distributes the traffic
leaving $A$, and the second stage forwards the traffic arriving at $B$ to $C$.
Let $q$ be the amount sent through $B$ and $f$ the fuel allocated to the
first stage. Their separate value functions are
\[
\begin{aligned}
U(p,q,f)&=\inf_{x_{AB},x_{AC}\ge0}
\left\{x_{AB}+4x_{AC}\;\middle|\;
\begin{aligned}
x_{AB}+x_{AC}&=p,\\
x_{AB}&=q,\\
x_{AB}+x_{AC}&\le f
\end{aligned}\right\},\\[1ex]
W(q,g)&=\inf_{x_{BC}\ge0}
\{2x_{BC}\mid x_{BC}=q,\ x_{BC}\le g\}.
\end{aligned}
\]
Taking $g=F-f$, the original problem is recovered by connecting the stages
and choosing their shared variables optimally:
\[
V(p,F)=\inf_{q,f\in\mathbb{R}}
\bigl(U(p,q,f)+W(q,F-f)\bigr).
\]
The costs add, and the internal choices $q$ and $f$ are eliminated by taking
an infimum. Indeed, the two fuel inequalities admit such an $f$ exactly when
$x_{AB}+x_{AC}+x_{BC}\le F$. All other constraints and costs are those of
the original problem. In this example the stage values are finite or
$+\infty$, so their sum does not require a convention for mixed infinities.

Although this composition is straightforward, its notation requires explicit
bookkeeping of interfaces, shared variables, and infima. Expanding $U$ and $W$
introduces the additional infima that define each stage.

We seek a language in which composition connects matching interfaces and the
semantics performs the minimization over shared variables. Only the external
interface remains exposed. This does not automatically compute the optimal
value; rather, it removes the need to write the minimization explicitly each
time two problems are composed. Such a language is
\textit{compositional by construction}.

Furthermore, as we will see in the next chapters, "eliminating" the inner
variables has a geometric interpretation. We keep the external values for which
some feasible choice of the inner variables exists. This operation is called
projection. In our example, this means keeping the pairs $(p,F)$ that admit
feasible choices of the flows and the shared variables $q,f$. For optimization,
we also need to keep track of costs; we will explain how after introducing epigraphs.
To reach that point, we need some formal definitions. The next section introduces
the necessary notions for us.

\section{Necessary notions}\label{sec:lp-notions}

We use two descriptions of the same class of value functions: extended
piecewise affine LP functions and right-hand side (RHS) parametric LP value
functions. We begin with the finite-valued piecewise affine case.

For the geometric arguments below, the \emph{epigraph} of a function
$f:\mathbb{R}^k\to\overline{\mathbb{R}}$ is
\[
\mathrm{epi}(f) := \{(p,t) \in \mathbb{R}^k \times \mathbb{R} \mid t \ge f(p)\},
\qquad \overline{\mathbb{R}}=\mathbb{R}\cup\{-\infty,+\infty\}.
\]
At a point where $f(p)=+\infty$, there is no real $t$ in the epigraph;
where $f(p)=-\infty$, every real $t$ is in it.

\begin{definition}[Piecewise Affine Convex Function]
A function $f : \mathbb{R}^n \to \mathbb{R}$ is called
\emph{piecewise affine convex} if there exist finitely many affine
functions $a_i^\top x + b_i$, $i = 1,\dots,k$, with $k\ge1$, such that
\[
f(x) = \max_{i=1,\dots,k} \{ a_i^\top x + b_i \}.
\]
\end{definition}

\begin{definition}[Extended Piecewise Affine LP Function]
An \emph{extended piecewise affine LP function} is a function
$f : \mathbb{R}^k \to \overline{\mathbb{R}}$, where
$\overline{\mathbb{R}} := \mathbb{R} \cup \{-\infty,+\infty\}$,
of one of the following forms:
\begin{enumerate}
\item For a nonempty polyhedron $D \subseteq \mathbb{R}^k$,
\[
f(p) = \begin{cases}
\max_{i=1,\dots,r}(a_i^\top p+\beta_i), & p \in D,\\
+\infty, & p \notin D,
\end{cases}
\qquad r \ge 1.
\]
\item For a nonempty polyhedron $D \subseteq \mathbb{R}^k$,
\[
f(p) = \begin{cases}
-\infty, & p \in D,\\
+\infty, & p \notin D.
\end{cases}
\]
\item $f(p)=+\infty$ for every $p$.
\end{enumerate}
The first form extends the preceding definition to a restricted domain;
the second and third forms account for unboundedness and infeasibility.
\end{definition}

\begin{definition}[RHS Parametric Linear Program]
A \emph{right-hand side parametric linear program} has value function
\[
v(p) := \inf_{x \in \mathbb{R}^n}
\{ c^\top x+d \mid Ax \ge Bp+e \},
\]
where $A \in \mathbb{R}^{m\times n}$, $B \in \mathbb{R}^{m\times k}$,
$c \in \mathbb{R}^n$, $e \in \mathbb{R}^m$, and $d \in \mathbb{R}$
are fixed. Thus the parameter appears only on the right-hand side.
We use $\inf\varnothing=+\infty$ for infeasibility and $v(p)=-\infty$
when the objective is unbounded below. The form $Ax\ge p$ is a special case.
\end{definition}

The following proposition shows that these definitions describe the same value
functions. This equivalence specifies the intended semantics of the graphical
language: extended piecewise affine LP functions represented by RHS parametric
linear programs.
\begin{proposition}
A function $v : \mathbb{R}^k \to \overline{\mathbb{R}}$ is an
extended piecewise affine LP function if and only if it is the value
function of a right-hand side parametric linear program as defined above.
\end{proposition}

\begin{proof}
\textbf{($\Rightarrow$)} Write the nonempty domain as
$D=\{p\mid Hp\ge h\}$.
For the first form, use the program
\[
v(p)=\inf_{t\in\mathbb{R}}
\{t\mid 0t\ge h-Hp,\quad
 t\ge a_i^\top p+\beta_i\text{ for }i=1,\dots,r\}.
\]
For $p\in D$, the minimum is the largest affine expression;
for $p\notin D$, the program is infeasible.
For the second form, minimize a free variable $z$ subject only to
$0z\ge h-Hp$. This has value $-\infty$ on $D$ and $+\infty$ outside it.
The third form is represented by an inconsistent constraint $0z\ge1$.
All these programs have the required affine RHS form.

\medskip
\textbf{($\Leftarrow$)} Let
$v(p)=\inf_x\{c^\top x+d\mid Ax\ge Bp+e\}$.
The feasible parameter domain
\[
D=\{p\mid \exists x:\ Ax\ge Bp+e\}
\]
is a polyhedron, since it is a projection of a polyhedron. Here projection
means keeping the parameter coordinates and eliminating the decision variables.
We use the standard facts that projections preserve polyhedra and that a
feasible LP with an objective bounded below attains a finite optimum
\cite{schrijver1986theory}.
If $D$ is empty, $v$ has the third form.
Otherwise, the recession cone of every feasible set is
$\{r\mid Ar\ge0\}$, independently of $p$. Its elements are directions $r$
such that moving from any feasible $x$ to $x+\lambda r$ remains feasible
for every $\lambda\ge0$.
By the standard recession criterion for linear programming
\cite{schrijver1986theory}, the objective
is unbounded below at a feasible parameter if and only if there exists
$r$ with $Ar\ge0$ and $c^\top r<0$.
Such a direction makes it unbounded below at every parameter in $D$,
giving the second form.

If no such direction exists, the optimum is finite and attained at every
$p\in D$. Consequently,
\[
\mathrm{epi}(v)=\{(p,t)\mid \exists x:\
Ax\ge Bp+e,\quad c^\top x+d\le t\}
\]
is a nonempty polyhedron. Write its defining inequalities as
$u_j^\top p+\alpha_j t\ge\gamma_j$.
Upward closure in $t$ implies $\alpha_j\ge0$.
The inequalities with $\alpha_j=0$ define the domain $D$;
those with $\alpha_j>0$ give finitely many affine lower bounds on $t$.
There is at least one such lower bound, since $v$ is finite on $D$.
Their maximum equals $v$ on $D$, giving the first form.
\end{proof}

These descriptions support different forms of reasoning. An RHS parametric
linear program specifies constraints and an objective, whereas an extended
piecewise affine LP function describes the geometry of its value function.
Different LP formulations may have the same value function. The intended
graphical language represents both descriptions and allows diagrams with the
same value function to be proved equal using its axioms.

The following proposition is the key technical lemma for the completeness
proof in the last section. It states that two parametric linear programs
with the same value function must share the same epigraph polyhedron.

\begin{proposition}\label{prop:equalepi}
Let
\[
v_i(y) := \inf_{x\in\mathbb{R}^{n_i}}\{c_i^\top x+d_i \mid A_i x \ge B_i y+e_i\},
\qquad i=1,2,
\]
where $A_i \in \mathbb{R}^{m_i \times n_i}$,
$B_i \in \mathbb{R}^{m_i \times k}$, $c_i\in\mathbb{R}^{n_i}$,
$e_i\in\mathbb{R}^{m_i}$, and $d_i\in\mathbb{R}$. If $v_1(y) = v_2(y)$ for all
$y \in \mathbb{R}^k$, then both programs admit a common representation
through the same epigraph polyhedron:
\[
v_i(y) = \inf\{t \mid (y,t) \in E\},
\qquad
E := \mathrm{epi}(v_1) = \mathrm{epi}(v_2).
\]
\end{proposition}

\begin{proof}
We use $\inf\varnothing=+\infty$ throughout.
At a feasible parameter, a finite optimum is attained; if the objective
is unbounded below, every real threshold $t$ admits a feasible point
with objective at most $t$. At an infeasible parameter, neither side
of the following equivalence holds. Thus, for each $i = 1,2$,
the epigraph of $v_i$ satisfies
\[
(y,t) \in \mathrm{epi}(v_i)
\iff
\exists x : A_i x \ge B_i y+e_i,\ c_i^\top x+d_i \le t,
\]
so $\mathrm{epi}(v_i)$ is the projection onto $(y,t)$ of the polyhedron
$P_i = \{(x,y,t) \mid A_i x \ge B_i y+e_i,\ c_i^\top x+d_i \le t\}$.
Projections of polyhedra are polyhedra, so each $\mathrm{epi}(v_i)$ is
a polyhedron. Since $v_1 = v_2$ pointwise, we have
$\mathrm{epi}(v_1) = \mathrm{epi}(v_2) =: E$.
Setting $E = \{(y,t) \mid Cy + at \ge h\}$ for suitable coefficients, we obtain
$v_i(y) = \inf\{t \mid Cy + at \ge h\}$ for both $i = 1, 2$.
\end{proof}

